% !TEX root = ../../../tree_notes.tex
% =====================================================================
%  13.3  Harvest at the peak: pair decomposition  sec:c10-harvest
%  Exports: prop:c10-pairtop  thm:c10-pairtop-count  lst:c10-peak
%           fig:c10-peaks  ex:c10-harvest-run
%           prob:c10-harvest-solved (CF 161D)  lst:c10-161d-soln
%  Ch.12's thm:c09-dpjoin (diameter arms joined at the peak) is the
%  special case; prop:c10-pairtop is the generalisation to arbitrary
%  additive path functionals.  Discharges Ch.16's forward-contract
%  reference to prop:c10-pairtop (ch11/99_chapter_end.tex L288).
% =====================================================================
\section{Path DP: Harvest at the Peak}\label{sec:c10-harvest}

\begin{previouslybox}
\begin{enumerate}[nosep, leftmargin=1.6em]
  \item Chapter~12's \thmref{c09-dpjoin} proved that the diameter of a tree
        can be computed by joining, at each node $v$, the two deepest
        downward arms through $v$. State the identity it uses for the
        length of a path through $v$. ($d(u,v)+d(v,w)=\mathit{arm}(u)+\mathit{arm}(w)$
        when $v$ is the peak of the path $u$--$w$)
  \item Fix a root $r$ and take any two distinct nodes $u, w$. Where is the
        unique node of the $u$--$w$ path closest to $r$? (their LCA
        $\ell$; every path has exactly one such \emph{peak})
  \item Chapter~2's \thmref{c01-height-recursion} counted a subtree by
        folding $\{1,\mathcal S(c)\}$; when the fold is over
        \emph{pairs of children} rather than children, what changes?
        (the answer becomes $\sum_{c_1<c_2}\mathit{sz}(c_1)\mathit{sz}(c_2)$
        plus a $\sum_c\mathit{sz}(c)$ for the pairs where one endpoint
        IS the peak -- the "arm of length zero" case)
  \item On the running tree $T$ (rooted at~$1$), where is the peak of the
        path $8$--$10$? ($\ell=1$; arms $8\to5\to2\to1$ (length~$3$) and
        $10\to9\to7\to1$ (length~$3$), total $d(8,10)=6$)
\end{enumerate}
\end{previouslybox}

\begin{topicintro}
Chapter~12 taught us that the diameter of a tree is computed by joining
two arms at a peak -- each node is a candidate peak, the two deepest
downward arms from it form a path through it, and the maximum over nodes
is the diameter. This section promotes that \emph{one} identity into a
\emph{scheme}: for \emph{any} functional $F$ on paths that decomposes as a
combination of its two arms around the peak,
\[
  \sum_{\{u,w\}\subseteq V} F(u,w)
  \;=\;
  \sum_{\ell\in V}\;\sum_{\substack{\{u,w\}: \mathrm{peak}(u,w)=\ell}} \mathrm{combine}\!\big(\mathit{arm}(u,\ell),\,\mathit{arm}(w,\ell)\big),
\]
and the inner sum factorises over the children of $\ell$ plus one
``arm of length $0$'' row for $\ell$ itself. This is \propref{c10-pairtop},
and it is the reason the whole ``count paths with property $P$'' family
on trees reduces to a single DFS. Ch.16's chapter-end cites this
proposition by name when it says ``'harvest at the LCA' was a proof
technique; from here it is a query technique'' -- the query machinery
there sits on top of the counting identity here.

\textbf{Prerequisites.} \thmref{c09-dpjoin} (Ch.12's diameter arms) as the
special case we generalise, not re-derive; \secref{c10-subtrees}'s
tuple-fold template; the depth/height/size tables of Ch.12's
\secref{c09-measures}. \emph{Not} assumed: LCA-as-a-data-structure
(Chapter~16 owns that); we use $\mathrm{peak}(u,w)$ as a \emph{definition},
and the DFS finds peaks implicitly via the parent guard.

\textbf{After this section you can}
\begin{itemize}[nosep]
  \item state \propref{c10-pairtop} and identify $\mathrm{combine}$ for a
        given problem (diameter, distance-$k$ count, sum-of-pair-distances);
        \hfill\lv{3}--\lv{4}
  \item cold-write \lstinline{dfs\_u\_returns\_hist}\ for CF~161D's
        distance-$k$ census in $\bigO{nk}$; \hfill\lv{4}
  \item explain exactly what the ``arm of length $0$'' row contributes
        and why forgetting it is the top WA cause; \hfill\lv{3}
\end{itemize}
\end{topicintro}

\subsection{Intuition}

Root the tree at any fixed $r$. Two facts become simultaneous.

\emph{Every unordered pair $\{u,w\}$ has a unique peak.} The $\mathrm{peak}(u,w)$ is
the unique node of the $u$--$w$ path closest to $r$; equivalently, it is
$\mathrm{LCA}(u,w)$. Uniqueness is a chain argument on the path: as you
travel from $u$ to $w$ you climb strictly toward $r$, reach the peak, then
descend strictly away. Two distinct minima would give two distinct
$u$--$w$ paths, impossible in a tree.

\emph{The path splits into two arms at the peak.} $\mathrm{arm}(u,\ell)$ is the
$\ell$-downward path to $u$; $\mathrm{arm}(w,\ell)$ similarly. Any
functional $F$ on $u$--$w$ that factors through these two arms -- length,
maximum, product-of-labels, colour-pattern -- satisfies
$F(u,w) = \mathrm{combine}(\mathrm{arm}(u,\ell),\mathrm{arm}(w,\ell))$.
Ch.12's \thmref{c09-dpjoin} chose $F(u,w)=d(u,w)$ and
$\mathrm{combine}(a,b)=\mathit{len}(a)+\mathit{len}(b)$, and maximised;
here we will \emph{sum} (or count, or $\max$) the same shape over every
pair, and the sum re-arranges as a sum over peaks, each peak harvesting
pairs of arms at itself.

The re-arrangement is the entire trick. A naive census is $\Theta(n^2)$
pairs; a peak census is $\Theta(n)$ peaks, each peak combining its
$\deg(\ell)$ arms. Since child sub-trees are disjoint, the arm-set at
$\ell$ is $\bigsqcup_{c\in\children(\ell)}\text{arms-at-}c$, plus the
$\ell$-itself arm of length~$0$. Merging children's arm-histograms one at a
time -- the same incremental trick as Ch.12's diameter join -- avoids the
$\binom{\deg}{2}$ explosion and gives $\bigO{k}$ (or $\bigO{\mathit{sz}}$)
per merge.

\subsection{Definitions and Notation}

\begin{definition}{Peak, arm, and arm functional}{c10-peak-def}
Fix a root $r$. For distinct nodes $u,w$, $\mathrm{peak}(u,w)$ is the
unique node on the $u$--$w$ path minimising $\depth(\cdot)$; equivalently
$\mathrm{peak}(u,w)=\mathrm{LCA}(u,w)$ in the rooted view. The
\term{arm} $\mathrm{arm}(u,\ell)$ is the $\ell$-downward simple path
$\ell=v_0,v_1,\dots,v_m=u$; we say it has \term{length} $m$ and \term{head}
$u$. We write $\mathrm{arm}(\ell,\ell)$ for the degenerate
arm of length~$0$ (the ``peak-as-endpoint'' case).
\end{definition}

\begin{proposition}{Harvest at the peak (pair decomposition of additive functionals)}{c10-pairtop}
Let $F$ be a functional on unordered pairs of distinct nodes of $T$ that
admits the arm-decomposition
$F(u,w)=\mathrm{combine}\!\big(\mathrm{arm}(u,\ell),\mathrm{arm}(w,\ell)\big)$
where $\ell=\mathrm{peak}(u,w)$ and $\mathrm{combine}$ is symmetric in its
two arguments and depends only on the two arms, not on the pair's identity.
Then
\[
  \sum_{\{u,w\}} F(u,w)
  \;=\;
  \sum_{\ell\in V}\;\sum_{\substack{\text{unordered }\{a,b\}\\ \text{arms at }\ell}} \mathrm{combine}(a,b),
\]
where the inner sum ranges over unordered pairs of distinct arms of $\ell$,
including the length-$0$ arm $(\ell,\ell)$ itself. \emph{The
generalisation of \thmref{c09-dpjoin}: the diameter is the $\max$-version,
this is the sum/count/$\max$-version.}
\end{proposition}
\begin{proofbox}
Two steps. (i)~\emph{Bijection between pairs and peak-arm-pairs}: the map
$\{u,w\}\mapsto(\mathrm{arm}(u,\ell),\mathrm{arm}(w,\ell))$ with
$\ell=\mathrm{peak}(u,w)$ is injective (arms determine $u$, $w$ and $\ell$)
and its image is exactly the set of unordered pairs of distinct arms of
$\ell$ across all $\ell$; no other pair of arms is realisable because two
arms of the same $\ell$ that share a first edge (same child subtree) do not
come from two nodes whose path peaks at $\ell$ -- their path peaks
\emph{below} $\ell$. So the peak of $\ell$ decomposes into
$\binom{\text{arms}}{2}=$ $\sum_{c_1<c_2}$ (across distinct children)
$+$ $\sum_{c}$ (the $(\ell,\ell)$ arm paired with each arm in $c$'s
subtree); this is why the length-$0$ arm must be present.
(ii)~\emph{Sum re-arrangement}: since $F$ is a function of the two arms
alone, $\sum_{\{u,w\}} F(u,w)$ equals $\sum_\ell \sum_{\text{arms at }\ell}\mathrm{combine}$
by (i). \qed
\end{proofbox}

\begin{theorem}{Counting pairs at distance $k$ via the peak}{c10-pairtop-count}
Given a tree $T$ and a target $k\ge 1$, let $H_v[j]$ denote the number of
nodes in $v$'s subtree at distance $j$ from $v$. Then the number of
unordered pairs $\{u,w\}$ with $d(u,w)=k$ is
\[
  \mathit{ans} \;=\; \sum_{v\in V}\;\sum_{\substack{c_1<c_2\\\mathstrut c_1,c_2\in\children(v)}}\;\sum_{i=0}^{k-1} H_{c_1}[i]\cdot H_{c_2}[k-2-i]
  \;+\;\sum_{v\in V}\;\sum_{c\in\children(v)} H_c[k-1],
\]
which is exactly \propref{c10-pairtop} instantiated to
$F(u,w)=[\,d(u,w)=k\,]$, $\mathrm{combine}(a,b) = [\,\mathit{len}(a)+\mathit{len}(b)=k\,]$,
and computable in $\bigO{nk}$ by \coderef{c10-peak}.
\end{theorem}
\begin{proofbox}
For a fixed peak $\ell$, an unordered pair of arms at $\ell$ of lengths
$i$ and $j$ contributes $[i+j=k]$ to the sum. Two structurally distinct
cases: (a)~both arms go into distinct child subtrees $c_1, c_2$; their
count of node-pairs is $H_{c_1}[i-1]\cdot H_{c_2}[j-1]$ (arms of length $i$
at $\ell$ correspond to nodes at distance $i-1$ from $c_1$), and we sum
over $i+j=k$; (b)~one arm is the degenerate length-$0$ arm at $\ell$ and
the other goes into a single child $c$, contributing $H_c[k-1]$. Cross-child
pairs with the same child ($c_1=c_2$) are excluded because their true peak
is \emph{below} $\ell$ -- those pairs are harvested at the LCA of the two
descendants, not at $\ell$; this is exactly the injectivity clause of
\propref{c10-pairtop}'s proof. Summing (a) and (b) over $\ell$ gives the
identity. \emph{Complexity}: each $H_v[\cdot]$ is a length-$(k+1)$
histogram; the incremental child-merge of \coderef{c10-peak} accumulates
$\bigO{k}$ per child, so $\bigO{k}$ per edge and $\bigO{nk}$ total. \qed
\end{proofbox}

\subsection{The Algorithm: One DFS, Arms Harvested at Every Peak}

\begin{lstlisting}[caption={Harvest at the peak: per-node arm histogram merged child-by-child; CF 161D's distance-k instance}, label={lst:c10-peak}]
// adj[u] = neighbours; call harvest(r, -1, k).  Returns nothing, writes ans.
vector<vector<int>> adj;
vector<vector<int>> H;                     // H[u][j] = nodes at distance j from u in u's subtree
long long ans;                             // LL: n*(n-1)/2 fits only in LL
int K;

void harvest(int u, int p) {
    H[u].assign(K + 1, 0);
    H[u][0] = 1;                           // u itself at distance 0
    for (int v : adj[u]) if (v != p) {
        harvest(v, u);                     // child histogram final first
        // harvest arms from u's already-merged side x arms from v's side
        // path length from a-side node (dist i from u) to b-side node
        // (dist j from v) equals i + j + 1 (the u-v edge is the +1).
        for (int i = 0; i <= K - 1; ++i)
            ans += (long long)H[u][i] * H[v][K - 1 - i];
        // merge v's histogram into u (shift by 1)
        for (int j = K; j >= 1; --j)
            H[u][j] += H[v][j - 1];
    }
}
// after harvest(r,-1,K): ans is the number of unordered pairs {x,y} with d(x,y)=K
\end{lstlisting}

\begin{complexitybox}
Time $\bigO{nk}$: each edge inspects $\bigO{k}$ histogram entries twice
(once for the cross-harvest, once for the shift-merge). Memory $\bigO{nk}$
for the histograms plus $\bigO{h}$ call stack; substitute \coderef{c09-iter}'s
two-phase form for chains at $n\ge 2\cdot 10^{5}$. Verdict: the distance-$k$
count is a $\bigO{nk}$ peak-harvest, and Ch.14 will prove this is optimal
for the histogram-merge family (their sum-of-subtree-size bound applies).
\end{complexitybox}

\subsection{Worked Example: Peak Census on $T$}

\begin{example}{Every unordered pair of $T$ grouped by peak and arm lengths}{c10-harvest-run}
Root $T$ at~$1$. For each $\ell\in V$, enumerate the arm-multiset
$\{\mathrm{arm}(u,\ell):u\in\text{subtree}(\ell)\}$ (with the length-$0$ arm at
$\ell$ included); a pair $\{u,w\}$ is harvested at $\ell=
\mathrm{peak}(u,w)$. The full census on $T$ has $\binom{10}{2}=45$ pairs
distributed across peaks:
\[
  \text{peak}=1:\; \textstyle\binom{4}{2}+\textstyle\binom{6}{1}\!\cdot\!1 = 6+6 = 12 \text{ pairs at peak }1\text{ from children }\{2,3\}\text{ plus arm }0.
\]
Rather than expand all $45$, we focus on the $k=2$ slice, which the
Ch.16 forward contract demands as \exref{c10-k2}. There are $\mathbf{10}$
pairs at distance~$2$ on $T$; their peaks are:
\end{example}

\begin{figure}[H]
\centering
\begin{tikzpicture}[diagram, scale=0.86, transform shape,
    tnode/.style={circle, draw=primaryblue, line width=0.9pt, fill=primaryblue!8,
                  minimum size=6.2mm, font=\small\bfseries, text=primaryblue!60!black},
    peaknode/.style={circle, draw=difforange, line width=1.4pt, fill=difforange!22,
                  minimum size=6.2mm, font=\small\bfseries, text=difforange!60!black}]
  \node[peaknode] (n1) at (0,0)     {$1$};
  \node[peaknode] (n2) at (-3,-1.6) {$2$};
  \node[peaknode] (n3) at ( 3,-1.6) {$3$};
  \node[tnode]    (n4) at (-5.2,-3.2) {$4$};
  \node[tnode]    (n5) at (-1.6,-3.2) {$5$};
  \node[tnode]    (n6) at ( 1.6,-3.2) {$6$};
  \node[peaknode] (n7) at ( 5.2,-3.2) {$7$};
  \node[tnode]    (n8) at (-1.6,-4.8) {$8$};
  \node[tnode]    (n9) at ( 5.2,-4.8) {$9$};
  \node[tnode]    (n10) at ( 5.2,-6.4) {$10$};
  \draw[tedge] (n1)--(n2) (n1)--(n3) (n2)--(n4) (n2)--(n5)
               (n3)--(n6) (n3)--(n7) (n5)--(n8) (n7)--(n9) (n9)--(n10);
  \node[annot, font=\scriptsize, text=difforange!70!black, anchor=west]
    at (7.4,0.2) {orange $=$ a peak $\ell$ that hosts some pair};
  \node[annot, font=\scriptsize, text=deepblue, anchor=west]
    at (7.4,-0.6) {peak~1: 5 pairs (1,4)(1,5)(1,6)(1,7)(2,3)};
  \node[annot, font=\scriptsize, text=deepblue, anchor=west]
    at (7.4,-1.4) {peak~2: 2 pairs (2,8)(4,5)};
  \node[annot, font=\scriptsize, text=deepblue, anchor=west]
    at (7.4,-2.2) {peak~3: 2 pairs (3,9)(6,7)};
  \node[annot, font=\scriptsize, text=deepblue, anchor=west]
    at (7.4,-3.0) {peak~7: 1 pair (7,10)};
  \node[annot, font=\scriptsize, text=warningred!70!black, anchor=west]
    at (7.4,-4.0) {total $5+2+2+1=10$ \checkmark};
\end{tikzpicture}
\caption{The $k=2$ census on $T$, harvested at each peak. The rows
$(1,4),(1,5),(1,6),(1,7),(2,8),(3,9),(7,10)$ each use the length-$0$ arm
at their peak (the peak itself is one endpoint); the remaining three
$(2,3),(4,5),(6,7)$ pair two genuine downward arms. The full \exref{c10-k2}
table lives in \secref{c10-distance}.}
\label{fig:c10-peaks}
\end{figure}

\begin{tracebox}[Fold order through \coderef{c10-peak} on $T$, root~$1$, $K=2$, ans starts at $0$]
post-order: 4, 8, 5, 2, 10, 9, 7, 6, 3, 1.  Histograms truncated to
length K+1=3.\\
Leaf 4: H[4]=(1,0,0).  Leaf 8: H[8]=(1,0,0).  Leaf 10: H[10]=(1,0,0).
Leaf 6: H[6]=(1,0,0).\\
5 (child 8): cross i=0..1: H[5][0]*H[8][1]+H[5][1]*H[8][0]=1*0+0*1=0.
Merge: H[5]=(1,1,0).\\
9 (child 10): cross 0. Merge: H[9]=(1,1,0).\\
7 (child 9): cross H[7][0]*H[9][1]+H[7][1]*H[9][0]=1*1+0*1=1 (pair (7,10)
with arm-length-0 on 7 side). Merge: H[7]=(1,1,1).\\
2 (children 4,5): after child 4: cross 1*0+0*1=0; H[2]=(1,1,0). After
child 5: cross H[2][0]*H[5][1]+H[2][1]*H[5][0]=1*1+1*1=2 (pairs (2,8)
arm-0 with dist 1 in child 5; and (4,5) with one arm in child 4 and one
in child 5). Merge: H[2]=(1,2,1).\\
3 (children 6,7): after child 6: cross 0; H[3]=(1,1,0). After child 7:
cross H[3][0]*H[7][1]+H[3][1]*H[7][0]=1*1+1*1=2 (pairs (3,9) arm-0; (6,7)
two-arm). Merge: H[3]=(1,2,1).\\
1 (children 2,3): after child 2: cross H[1][0]*H[2][1]=1*2=2 (pairs
(1,4),(1,5)). Merge H[1]=(1,1,2). After child 3: cross
H[1][0]*H[3][1]+H[1][1]*H[3][0]=1*2+1*1=3 (pairs (1,6),(1,7),(2,3)).\\
Total ans = $1$ (at~7) $+\,2$ (at~2) $+\,2$ (at~3) $+\,2$ (at~1, first child)
$+\,3$ (at~1, second child) $= \mathbf{10}$. \checkmark\ matches
\figref{c10-peaks} census and Ch.16's forward reference.
\end{tracebox}

\begin{insightbox}[Key takeaways]
\begin{enumerate}[nosep, leftmargin=1.6em]
  \item \propref{c10-pairtop} is Ch.12's \thmref{c09-dpjoin} with $\max$
        replaced by $\mathrm{combine}$; the two-arm-at-the-peak geometry
        is \emph{the same}.
  \item The length-$0$ arm at $\ell$ (the ``one endpoint IS the peak''
        case) is what separates \coderef{c10-peak}'s cross-product from a
        plain $\sum_{c_1<c_2}$; forgetting it drops $n-1$ contributions on
        a chain and is the top WA cause.
  \item The incremental merge (harvest $u$'s already-accumulated side
        against $v$'s side, \emph{then} shift-merge $v$ into $u$)
        auto-handles the ``arms in the same child subtree'' exclusion:
        those pairs are simply never cross-multiplied.
\end{enumerate}
\end{insightbox}

\subsection{Variations and Twists}

\begin{itemize}[nosep]
  \item \textbf{Sum of pair distances} ($F(u,w)=d(u,w)$): set
        $\mathrm{combine}(a,b)=\mathit{len}(a)+\mathit{len}(b)$ in
        \propref{c10-pairtop}. The double sum telescopes into a per-edge
        contribution $\sum_{e} \mathit{sz}_e\,(n-\mathit{sz}_e)$: one DFS,
        no histogram. \hfill\lv{3}
  \item \textbf{Count pairs with peak exactly at $v$}: the answer at $v$
        is $\binom{\mathit{sz}(v)}{2}-\sum_c\binom{\mathit{sz}(c)}{2}$;
        scalar fold. \hfill\lv{3}
  \item \textbf{Diameter-as-count}: instantiate $F$ to
        $[d(u,w)=D]$ to get Ch.12's diameter census; \thmref{c09-dpjoin}
        recovers it. This is the ``same proof, different combine''
        discipline. \hfill\lv{3}
  \item \textbf{Centroid decomposition upgrade}: when $k$ is
        \emph{unbounded} and $n\ge 10^{5}$, the $\bigO{nk}$ budget
        explodes; Ch.24 (\chref{c18}) shows how to answer distance
        queries in $\bigO{n\log n}$ via centroid recursion. Ch.13 does
        not go there. \hfill\lv{4}
\end{itemize}

\subsection{Common Pitfalls}

\begin{warningbox}[Trap 1 -- the missing length-$0$ arm at the peak]
Writing \lstinline{for i in 0..K-1: ans += H[u][i] * H[v][K-1-i]} at the
top of the child loop and \lstinline{H[u][0]=1} before the loop handles
both cases correctly -- but \emph{only} if \lstinline{H[u][0]=1} is set
BEFORE any child is merged. A common variant of the bug: initialising
$\mathit{H}[u][0]=0$ and setting it after the loop; then the
arm-at-$u$-itself contribution vanishes. Symptom: every ``one endpoint
is the peak'' pair (the degenerate row in \figref{c10-peaks}) is missing
from the answer; on a chain this drops $n-1$ pairs, on $T$ at $k=2$ it
drops $7$ out of $10$.
\end{warningbox}

\begin{warningbox}[Trap 2 -- cross-multiplying arms in the same child]
Merging all children's histograms first and then taking a single
$\sum_{i+j=k}\mathit{H}[u][i]\mathit{H}[u][j]$ counts, for each pair of
descendants of the same child $c$, \emph{twice} at $u$ (once for each
order) and misses their true peak (which is below $u$). Correct: harvest
\emph{before} merging, so at the moment of harvesting $c_i$'s contribution
only \emph{earlier-merged} children's arms are visible at $u$. This is the
incremental trick of \coderef{c10-peak}; Ch.12's \thmref{c09-dpjoin} uses
the same discipline for the $\max$.
\end{warningbox}

\begin{warningbox}[Trap 3 -- $\mathit{int}$ overflow on the pair count]
$n=10^{5}$ makes $\binom{n}{2}\approx 5\cdot 10^{9}$, past $\mathit{int}$.
\lstinline{ans} must be \lstinline{long long}, and each per-merge product
$\mathit{H}[u][i]\cdot\mathit{H}[v][k-1-i]$ must be cast before
multiplication (two \lstinline{int}s whose product needs 64-bit
\texttt{long\,long} but overflows \lstinline{int} silently). Symptom: sample small,
submits big, WA on hidden with $\mathit{ans}$ looking ``about half what
it should be'' -- the upper 32 bits dropped.
\end{warningbox}

\subsection{Problem Lab}

The section's anchor is Codeforces~161D, \probref{c10-harvest-solved}: the
distance-$k$ pair count on a tree. It is \thmref{c10-pairtop-count}
straight off the page, and the peak-harvest in \figref{c10-peaks} is
exactly the algorithm's run on $T$.

% ---------------------------------------------------------------------------
\solvedproblem{c10-harvest-solved}{Distance in Tree}{Codeforces 161D}
  {https://codeforces.com/problemset/problem/161/D}{CF 1800 L4}
  {harvest at the peak, arm histograms, incremental child merge, long long}

\begin{problemstatement}[Distance in Tree (samples verified against the live problem)]
Given a tree with $n$ nodes and an integer $k\ge 1$, count the number of
unordered pairs of distinct nodes $(u,v)$ with $d(u,v)=k$ (edges
counted). Constraints: $1\le n\le 5\cdot 10^{4}$, $1\le k\le 500$. Input:
$n$ and $k$, then $n-1$ lines of edges. Output: a single integer.
\textbf{Sample~1}: $n=5$, $k=2$, edges $1\!-\!2,2\!-\!3,3\!-\!4,2\!-\!5$
$\Rightarrow$ output \texttt{4} (pairs $(1,3),(1,5),(2,4),(3,5)$).
\textbf{Sample~2}: $n=5$, $k=3$, edges $1\!-\!2,2\!-\!3,3\!-\!4,4\!-\!5$
$\Rightarrow$ output \texttt{2} (pairs $(1,4),(2,5)$).
\end{problemstatement}

\begin{attemptgate}[Attempt it yourself first -- \lv{4}, give yourself 25 minutes]
Do not jump to $\Theta(n^2)$ pair enumeration. Two questions unlock
$\bigO{nk}$: (i) at which node is a pair harvested? (ii) what does the
state at that node need to remember about each arm?
\peek{Given a pair $\{u,w\}$ with $d(u,w)=k$, where on the path
      $u$--$w$ do you naturally want to ``count'' the pair exactly once?}
\nudge{At $\ell=\mathrm{peak}(u,w)$. Then $u$ and $w$ sit in two
      \emph{distinct} downward arms of $\ell$, or one of them IS $\ell$.}
\fullhint{Keep a histogram $\mathit{H}[v][j]=$ number of nodes at
      distance $j$ from $v$ in $v$'s subtree. Post-order DFS: at $u$,
      for each child $v$, first harvest the cross product
      $\sum_i\mathit{H}[u][i]\cdot\mathit{H}[v][k-1-i]$ (this counts
      pairs with one endpoint already merged into $u$'s side and the other
      in $v$'s side; the $+1$ from the $uv$ edge is folded into the
      histogram shift), \emph{then} shift-merge $v$'s histogram into $u$:
      $\mathit{H}[u][j]\mathrel{+}= \mathit{H}[v][j-1]$. Initialise
      $\mathit{H}[u][0]=1$ at every node \emph{before} the child loop so
      the ``one endpoint is $u$'' row is counted automatically. Answer
      lives in a global \lstinline{long long} \lstinline{ans}. Total
      $\bigO{nk}$.}
\end{attemptgate}

\solstep{Understanding the Problem}
The naive $\Theta(n^2)$ pair enumeration hits $n=5\cdot 10^{4}$ at
$\approx 1.25\cdot 10^{9}$, borderline TLE. Two structural facts collapse
it. (i) Every pair has a unique peak, so $\sum_{\text{pairs}}$ can be
re-grouped as $\sum_{\text{peaks }\ell}\sum_{\text{arms at }\ell}$ by
\propref{c10-pairtop}. (ii) At $\ell$, the only information needed per
arm is its \emph{length}, and there are $k$ relevant lengths; that is
$\bigO{k}$ per peak. Total: $\bigO{nk}$.

\solstep{Key Observations}
\begin{observation}
$F(u,w) = [\,d(u,w)=k\,]$ factors as $\mathrm{combine}(a,b)=[\,\mathit{len}(a)+\mathit{len}(b)=k\,]$,
so \propref{c10-pairtop} applies verbatim.
\end{observation}
\begin{observation}
The child-subtree histograms are the right state: they carry $\bigO{k}$
numbers per node, they shift by $+1$ at every parent edge, and they
combine via a cross-product. The cross-product is $\bigO{k}$ per child.
\end{observation}
\begin{observation}
Incremental merging (harvest \emph{before} shifting the next child in)
avoids the $\binom{\deg}{2}$ pair-of-children explosion, and it makes the
length-$0$ arm at $u$ contribute automatically through $\mathit{H}[u][0]=1$.
\end{observation}

\solstep{Think It Through With Me}
\begin{thinkbox}
\deadend{Enumerate all $\binom{n}{2}$ pairs, compute each distance via
      LCA-based $\bigO{1}$ queries, count.}
{Requires Chapter~16's LCA pre-processing (this chapter does not assume
it) and costs $\binom{n}{2}\approx 1.25\cdot 10^{9}$ -- well past the
$2$-second CF limit even if the LCA is $\bigO{1}$. Also misses the
point: the peak structure is what makes the problem tractable, and
pair enumeration discards it.}
\deadend{DFS once, then for each node count pairs at distance $k$ with
      one endpoint AT the node by looking at $\mathit{H}[u][k]$.}
{That counts each pair \emph{twice} (once at each endpoint), giving
$\mathit{ans}/2$, and only counts pairs with one endpoint an ancestor of
the other -- the pairs with a strictly-below peak are missing entirely.}
\aha{Harvest at the peak via \thmref{c10-pairtop-count}: one histogram
     per node, cross-product on entry of each child, then shift-merge.
     $\bigO{nk}$ and clean.}
\end{thinkbox}

\solstep{Solution}
Post-order DFS. At $u$: set $\mathit{H}[u][0]=1$. For each child $v$, first
add to $\mathit{ans}$ the harvest $\sum_{i=0}^{k-1}\mathit{H}[u][i]\cdot\mathit{H}[v][k-1-i]$,
then shift-merge $\mathit{H}[v][j-1]\to\mathit{H}[u][j]$ for
$j=k,\dots,1$. Answer is $\mathit{ans}$. Chain-safe by substituting
\coderef{c09-iter}'s two-phase skeleton for the recursion.

\solstep{Proof of Correctness}
\begin{proofbox}[\coderef{c10-peak}'s output equals the number of pairs at distance $k$]
Follows \thmref{c10-pairtop-count} with the incremental-merge
justification. At each node $u$, we claim: \emph{before processing child
$v$}, $\mathit{H}[u][i]$ is the number of nodes in $u$'s ``already-merged
side'' at distance $i$ from $u$ (which is $\{u\}$ union the subtrees of
all previously-processed children, with distances measured correctly).
\emph{Base}: before the first child, already-merged side $=\{u\}$ and
$\mathit{H}[u][0]=1$. \emph{Induction}: shift-merge adds $v$'s subtree
with distance $+1$ (a node at distance $j-1$ from $v$ is at distance $j$
from $u$ via edge $uv$). At the moment of harvest, pairs counted are
exactly those with one endpoint in the already-merged side and the other
in $v$'s subtree, each at its correct distance from $u$; their sum of
distances from $u$ is $i + (k-1-i) + 1 = k$ (the $+1$ accounts for the
$uv$ edge on the $v$-side arm; the $u$-side arm already counts distance
\emph{from $u$} directly). No pair is counted twice because each is
harvested exactly at its peak, and each peak processes each child exactly
once. Pairs with one endpoint \emph{equal to} $u$ are covered by the
$i=0$ term ($\mathit{H}[u][0]=1$ and $\mathit{H}[v][k-1]$). \qed
\end{proofbox}

\solstep{Implementation}
\begin{lstlisting}[caption={CF 161D -- harvest-at-peak with per-node histogram, iterative to survive $n=5\cdot 10^{4}$ chains}, label={lst:c10-161d-soln}]
#include <bits/stdc++.h>
using namespace std;
using ll = long long;

int main() {
    int n, K;  cin >> n >> K;
    vector<vector<int>> adj(n);
    for (int i = 0; i < n - 1; ++i) {
        int u, v;  cin >> u >> v;  --u; --v;
        adj[u].push_back(v);  adj[v].push_back(u);
    }
    // two-phase iterative DFS to get a safe fold order (chain-proof)
    vector<int> par(n, -1), order;  order.reserve(n);
    vector<int> stk{0};
    while (!stk.empty()) {
        int u = stk.back();  stk.pop_back();
        order.push_back(u);
        for (int v : adj[u]) if (v != par[u]) { par[v] = u; stk.push_back(v); }
    }
    vector<vector<int>> H(n, vector<int>(K + 1, 0));
    ll ans = 0;
    for (int idx = (int)order.size() - 1; idx >= 0; --idx) {   // reversed preorder = post-order on a tree
        int u = order[idx];
        H[u][0] = 1;                                            // the length-0 arm at u
        for (int v : adj[u]) if (v != par[u]) {
            // harvest u's already-merged side x child v's side
            for (int i = 0; i <= K - 1; ++i)
                ans += (ll)H[u][i] * H[v][K - 1 - i];
            // shift-merge: v's distance-(j-1) nodes are u's distance-j nodes
            for (int j = K; j >= 1; --j)
                H[u][j] += H[v][j - 1];
        }
    }
    cout << ans << '\n';
}
\end{lstlisting}

\solstep{Trace}
\begin{tracebox}[Sample~1: $n=5$, $k=2$, edges $1\!-\!2,2\!-\!3,3\!-\!4,2\!-\!5$ $\Rightarrow$ 4]
Root at~1, children: 1$\to$$\{2\}$, 2$\to$$\{3,5\}$, 3$\to$$\{4\}$.
Post-order: 4, 3, 5, 2, 1. Histograms truncated to length 3.
Node 4 (leaf): H[4]=(1,0,0).
Node 3 (child 4): cross H[3][0]*H[4][1]=1*0=0. Merge H[3]=(1,1,0).
Node 5 (leaf): H[5]=(1,0,0).
Node 2 (children 3, 5): after child 3: cross H[2][0]*H[3][1] + H[2][1]*H[3][0]
  = 1*1 + 0*1 = 1 (pair (2,4)). Merge H[2]=(1,1,1). After child 5:
  cross H[2][0]*H[5][1] + H[2][1]*H[5][0] = 1*0 + 1*1 = 1 (pair (3,5)
  via arms at 2 of lengths 1 and 1). Merge H[2]=(1,2,1).
Node 1 (child 2): cross H[1][0]*H[2][1] + H[1][1]*H[2][0] = 1*2 + 0*1 = 2
  (pairs (1,3) and (1,5), both use the length-0 arm at 1). Merge.
Total ans = $1$ (at~2, with child~3) $+\,1$ (at~2, with child~5)
$+\,2$ (at~1) $= \mathbf{4}$. \checkmark\ matches official.
\end{tracebox}

\begin{tracebox}[Sample~2: $n=5$, $k=3$, edges $1\!-\!2,2\!-\!3,3\!-\!4,4\!-\!5$ $\Rightarrow$ 2]
Chain rooted at~1. Children: 1$\to$$\{2\}$, 2$\to$$\{3\}$, 3$\to$$\{4\}$,
4$\to$$\{5\}$. Post-order: 5, 4, 3, 2, 1. Histograms truncated to length 4.
Node 5: H=(1,0,0,0).
Node 4 (child 5): cross $\sum_{i=0}^{2}$ H[4][i]*H[5][2-i] = all zeros
  (H[5][j]=0 for $j\ge 1$). Merge H[4]=(1,1,0,0).
Node 3 (child 4): cross $\sum_{i=0}^{2}$ H[3][i]*H[4][2-i] = H[3][0]*H[4][2]
  + H[3][1]*H[4][1] + H[3][2]*H[4][0] = 1*0 + 0*1 + 0*1 = 0. Merge
  H[3]=(1,1,1,0).
Node 2 (child 3): cross H[2][0]*H[3][2] + H[2][1]*H[3][1] + H[2][2]*H[3][0]
  = 1*1 + 0*1 + 0*1 = 1 (pair (2,5), arms lengths 0 and 3 at peak 2,
  i.e. one endpoint IS the peak). Merge H[2]=(1,1,1,1).
Node 1 (child 2): cross H[1][0]*H[2][2] + H[1][1]*H[2][1] + H[1][2]*H[2][0]
  = 1*1 + 0*1 + 0*1 = 1 (pair (1,4), arm-0 at peak 1). Merge.
Total = $1 + 1 = \mathbf{2}$. \checkmark\ matches official.
\end{tracebox}

\solstep{Complexity}
\begin{complexitybox}
Time $\bigO{nk}$ ($n\le 5\cdot 10^{4}$, $k\le 500$: $\le 2.5\cdot 10^{7}$,
comfortable). Memory: $\bigO{nk}$ for histograms ($\le 2.5\cdot 10^{7}$
\lstinline{int}s $\approx 100$ MB -- if CF's limit is 256 MB, fits; swap
\lstinline{int} for \lstinline{short} if you want headroom since
$\mathit{H}[u][j]\le n$ but actually $\mathit{H}[u][j]\le$ subtree size
$\le n$ and $n=5\cdot 10^{4}$ needs 16 bits, borderline; use
\lstinline{int} here). Verdict: \ac. The peak structure is the whole
solution; \thmref{c10-pairtop-count} is the one-line justification.
\end{complexitybox}

\solstep{Pitfalls and Variants}
\begin{warningbox}
Three ways to lose this problem: (i) merge child \emph{before} harvesting,
which loses the ``cross-arms'' pairs and double-counts within-child pairs
(Trap~2 above); (ii) $\mathit{H}[u][0]=1$ after the child loop, which
drops the peak-as-endpoint rows (Trap~1); (iii) $\mathit{int}$ \lstinline{ans}
-- the count reaches $n-1$ on a chain at $k=1$, and $n^2/4$ on a star at
$k=2$, both past $\mathit{int}$ at $n=5\cdot 10^{4}$. Variant worth doing
after this: CSES~2081 ``Fixed-Length Paths II'' (Ch.12 owns the URL) is
the same harvest-at-peak identity with $F(u,w)=[\,a\le d(u,w)\le b\,]$;
the histogram merge is unchanged, only the $\sum$'s bounds shift.
\end{warningbox}

\subsection{Practice Problems}

\begin{stuckbox}
\begin{itemize}[nosep, leftmargin=1.4em]
  \item Is your child loop \emph{harvest then merge}, not \emph{merge
        then harvest}? The order is what keeps the cross-arms invariant.
  \item Does $\mathit{H}[u][0]=1$ appear BEFORE the child loop? After
        the loop is a Trap~1 WA.
  \item Is the harvest loop \lstinline{i in [0, K-1]} (inclusive) with
        \lstinline{H[v][K-1-i]}? A common off-by-one writes
        \lstinline{H[v][K-i]}, which effectively asks $i+j+0=k$ (forgetting
        the $uv$ edge) and silently drops pairs.
\end{itemize}
\end{stuckbox}

\begin{practiceset}
  \practice{Child and Puzzle}{Codeforces 839B}
    {https://codeforces.com/problemset/problem/839/B}{CF 1700 L3}
    {Same peak-decomposition idea, scalar version. For each node, you need
    the \emph{top two downward arm lengths}; combining them at the peak
    gives Ch.12's diameter, but the problem asks whether the top-two
    lengths at every node satisfy a per-node budget. Do the fold as a
    three-coordinate tuple $\mathcal S(v)=(\mathit{top}_1,\mathit{top}_2,
    \mathit{ok})$ -- it is \propref{c10-pairtop} with $F$ reduced to a
    predicate on top-two arms.}
\end{practiceset}

\paragraph{Back-cites (URLs live elsewhere in the book).}
\begin{itemize}[nosep,leftmargin=1.4em]
  \item \textbf{CSES 1130 ``Tree Diameter''}: Ch.12 owns the pointer in its
        mixed set. Re-solve it now that you see \thmref{c09-dpjoin} as a
        special case of \propref{c10-pairtop}: $\mathrm{combine}(a,b)=
        \mathit{len}(a)+\mathit{len}(b)$, $\max$ over peaks.
  \item \textbf{LC 1245 ``Tree Diameter''} (\chref{c09} owns the URL): same
        reduction as CSES~1130 -- Ch.12's \thmref{c09-twosweep} is the
        two-BFS algorithm; you now know the DP one and can see they agree.
\end{itemize}

\begin{recallbox}[10-second refresher -- 13.3 Harvest at the Peak]
Every pair has a unique peak $\ell=\mathrm{peak}(u,w)=\mathrm{LCA}(u,w)$,
and any additive path functional $F$ factors as
$\mathrm{combine}(\mathrm{arm}(u,\ell),\mathrm{arm}(w,\ell))$
(\propref{c10-pairtop}, Ch.12's \thmref{c09-dpjoin} generalised).
$\sum_{\{u,w\}}F(u,w)=\sum_\ell\sum_{\text{arms at }\ell}\mathrm{combine}$.
Distance-$k$ count instantiates $\mathrm{combine}=[i+j=k]$ and takes
$\bigO{nk}$ via \coderef{c10-peak}: per-node histogram, harvest BEFORE
merge, $\mathit{H}[u][0]=1$ for the length-$0$ arm. CF~161D
(\probref{c10-harvest-solved}) is the canonical problem. The
length-$0$-arm row is what Ch.16's \exref{c10-k2} forward reference will
flag as ``degenerate''.
\selftest{
  \item State \propref{c10-pairtop} in one sentence and name the Ch.12
        theorem it generalises.
  \item What are the two structural cases at peak $\ell$ in the distance-$k$
        census? (Cross-child arms and one arm equals $\ell$ itself.)
  \item Why does the incremental ``harvest-then-merge'' order exclude
        same-child pairs at $\ell$?
  \item Which row of \figref{c10-peaks} would disappear if you set
        $\mathit{H}[u][0]=0$?
  \item Give the $\bigO{}$ of \coderef{c10-peak} in $n$ and $k$ and name
        the constraint signature it fits.
}
\end{recallbox}
